% Example XeLaTeX document configured for Greek (UTF-8)
% Compile with: xelatex report.tex  (or use latexmk -xelatex)
\documentclass[12pt,a4paper]{article}
\usepackage{fontspec}
\usepackage{polyglossia}
\setdefaultlanguage{greek}
\setotherlanguage{english}

% Choose a font with good Greek coverage. Noto Serif is a good default if installed.
\newfontfamily\greekfont[Script=Greek]{Noto Serif}
\newfontfamily\englishfont{DejaVu Sans}
% Define monospace font that contains Greek glyphs to avoid polyglossia errors
\newfontfamily\greekfonttt[Script=Greek]{DejaVu Sans Mono}

\usepackage{microtype}
\usepackage{geometry}
\geometry{margin=2.5cm}

\title{Δείγμα Αναφοράς}
\author{Όνομα Φοιτητή}
\date{\today}

\begin{document}
\maketitle

\section{Εισαγωγή}
Αυτό είναι ένα μικρό παράδειγμα εγγράφου στα Ελληνικά. Χρησιμοποιούμε XeLaTeX (ή LuaLaTeX) και το πακέτο \texttt{polyglossia} για σωστή υποστήριξη της Ελληνικής γλώσσας και Unicode.

\section{Παράδειγμα Μαθηματικών}
Μπορούμε να γράψουμε μαθηματικά κανονικά:
\[ E = mc^2 \]

\section{Λίστα}
\begin{itemize}
  \item Πρώτο σημείο
  \item Δεύτερο σημείο
\end{itemize}

\end{document}
